% !TEX program = tectonic
% !TEX root = forschungsbericht.tex
\documentclass[type=forschungsbericht]{zukunftbau}

\usepackage{xurl}
% LR-019: biblatex justifies bibliography entries; its \biburlsetup sets
% \Urlmuskip=0mu plus 3mu, so every URL breakpoint stretches to fill the line
% and URLs render with visible gaps ("www . baunetzwissen . de"). Reset the
% muskip after biblatex's setup so URLs stay continuous (xurl still allows the
% breaks, just without stretch).
\usepackage{etoolbox}
\appto\biburlsetup{\Urlmuskip=0mu\relax}
\addbibresource{references.bib}

% LR-020: this is a digital deliverable; drop the print-signature padding so
% the PDF ends right after the bibliography instead of on trailing blank pages.
\SemioSignaturePadding{false}
% Digital deliverable: no recto-forcing blank versos before appendix chapters.
\SemioOpenAny{true}

\title{Entwerfen mit Bestand}
\subtitle{Eine offene Plattform für einen KI-unterstützten, performance-optimierten und integrativen Entwurfsprozess mit wiederverwendeten Baukomponenten}
\antragstellendeinstitution{%
  \AddressOrg{%
    Leibniz Universität Hannover\\
    Fakultät für Architektur und Landschaft\\
    Institut für Entwerfen und Konstruieren\\
    Abteilung Nachhaltige Gebäudesysteme%
  }%
  \AddressPeople{%
    Institutsleitung: Prof. Dr.-Ing. Philipp Geyer\\
    Projektleitung: Ueli Saluz\\
    Weitere Bearbeitung: Nikolaus Möllenhof%
  }%
}
\kooperationspartner{%
  \AddressOrg{%
    Universität der Künste Berlin\\
    Fakultät für Architektur\\
    Institut für Architektur und Städtebau\\
    Konstruktives Entwerfen und Tragwerksplanung%
  }%
  \AddressPeople{%
    Institutsleitung: Prof. Dr.-Ing. Christoph Gengnagel\\
    Bearbeitung: Kinan Sarakbi%
  }%
}
\aktenzeichen{10.08.18.7-25.06}
\foerderzeitraum{03.11.2025-02.05.2027}
% Anlage 8 macht eine DOI bzw. einen anderen Persistent Identifier zur Pflicht.
% Platzhalter bis zur Vergabe — bewusst sichtbar, nicht leer.
\doi{10.0000/zzz-doi-ausstehend}
\footerauthors{Ueli Saluz, Kinan Sarakbi}
\footerlogos{asset/logo/luh.png, asset/logo/udk.png}
\makeatletter
\renewcommand{\semio@cover@zukunftbau@logo}{asset/logo/zukunft-bau.png}
\renewcommand{\semio@cover@bbsr@logo}{asset/logo/bbsr.png}
\renewcommand{\semio@cover@bundesministerium@logo}{asset/logo/gefördert-durch-bmwsb.png}
\renewcommand{\semio@cover@arbeitsprobe@image}{asset/demonstrator/aggegator.png}
\makeatother
% Deckblatt-Blurb, nicht der Abstract: das Beschreibungsfenster des Deckblatts
% ist eine schmale Zeile im Fensterraster. Die vollständige Kurzbeschreibung und
% der englische Abstract stehen dort, wo Anlage 8 sie verlangt — im gleichnamigen
% Kapitel.
\kurzfassung{%
  Eine offene End-To-End-Plattform, die Bauteilkatalog, Entwurfsumgebung und Leistungsbewertung verbindet und den Entwurf mit wiederverwendeten Baukomponenten von der Datenaufnahme bis zur Vorbewertung durchgängig unterstützt.%
}

\ExplSyntaxOn
% LR-2.16: runs under \ExplSyntaxOn where the space char is catcode 9 (ignored),
% so plain spaces are eaten and the header renders "VerwendungimBericht". Use
% explicit ~ space tokens so the words stay separated.
\cs_set:Npn \semio_glossary_header_definition: { Verwendung~im~Bericht }
\ExplSyntaxOff

\GlossaryDefine{Baukomponente}{Physisches Bauteil oder zusammengehörige Gruppe von Bauteilen, die im Entwurf als eine Einheit betrachtet und wiederverwendet wird, etwa eine Stütze, eine Decke oder ein Rahmenelement.}
\GlossaryDefine{Entwurfsrolle}{Funktion, die eine Baukomponente in einem neuen Entwurf übernehmen soll. Sie kann von der ursprünglichen Verwendung abweichen, etwa wenn ein Träger als Stütze eingesetzt wird.}
\GlossaryDefine{Bestandsquelle}{Ort oder Kanal, aus dem wiederverwendbare Baukomponenten stammen, etwa eine Bauteilbörse, ein Rückbauinventar oder eine direkte Bestandsaufnahme vor Ort.}
\GlossaryDefine{Bauteilbörse}{Organisation, Lager oder digitaler Angebotskanal, über den wiederverwendbare Baukomponenten dokumentiert, vermittelt oder verkauft werden.}
\GlossaryDefine{Qualifizierung}{Schrittweise Erfassung, Prüfung und nachvollziehbare Ergänzung von Bauteilinformationen, bis sie für Entwurf und Bewertung nutzbar sind.}
\GlossaryDefine{Typologiefamilie}{Ordnung von Baukomponenten nach ihrer geometrisch-konstruktiven Grundform: lineare und flächige Elemente, Knoten- und Übergangselemente sowie Rahmen- und Systemelemente.}
\GlossaryDefine{Nachweiszustände}{Kennzeichnung, wie belastbar eine einzelne Angabe ist: vorhanden, abgeleitet, angenommen oder noch offen.}
\GlossaryDefine{Leistungsbewertung}{Einschätzung ausgewählter Eigenschaften oder Auswirkungen einer Entwurfsvariante auf Grundlage der verfügbaren Informationen, etwa zu Energie oder Tragwerk.}
\GlossaryDefine{Vorbewertung}{Frühe, vereinfachte Leistungsbewertung, die im Entwurf der Orientierung dient und keine vollständige fachliche Prüfung ersetzt.}
\GlossaryDefine{Bauteilobjekt}{Digitales Modell einer Baukomponente, das Geometrie und zugehörige fachliche Angaben in einer gemeinsamen Struktur bündelt.}
\GlossaryDefine{Generator}{Automatisches Verfahren, das aus erfassten Bauteilangaben ein digitales Bauteilobjekt mit den benötigten geometrischen Darstellungen erzeugt.}
\GlossaryDefine{Einspeiseplattform}{Digitaler Eingangsbereich der Plattform aus Bauteilportal und Wizard, in dem Bestandsdaten erfasst, übernommen und aufbereitet werden.}
\GlossaryDefine{Entwurfswerkzeug}{Digitaler Arbeitsbereich aus Bauteilkatalog, Entwurfsumgebung und Konfigurator zur Auswahl, Kombination und Bewertung von Baukomponenten.}
\GlossaryDefine{Bauteilportal}{Weboberfläche zur manuellen Erfassung, Übernahme und Qualifizierung von Bestandsdaten als Grundlage für digitale Bauteilobjekte.}
\GlossaryDefine{Wizard}{Benutzeroberfläche für einen geleiteten Abfrageprozess, mit der unvollständige oder unstrukturierte Bauteildaten Schritt für Schritt erfasst und ergänzt werden.}
\GlossaryDefine{Bauteilkatalog}{Übersichts- und Suchbereich des Entwurfswerkzeugs, in dem verfügbare Baukomponenten und ihre Angaben gefunden und ausgewählt werden.}
\GlossaryDefine{Entwurfsumgebung}{Arbeitsbereich des Entwurfswerkzeugs, in dem Baukomponenten räumlich angeordnet, kombiniert und zu Entwurfsvarianten weiterentwickelt werden.}
\GlossaryDefine{Konfigurator}{Austauschbereich des Entwurfswerkzeugs für den Kontakt mit Bauteilbörsen, etwa für Anfragen, Reservierungen, Statusänderungen und projektspezifische Angaben.}
\GlossaryDefine{semio}{Quelloffenes Software-Toolkit zur Beschreibung, Verknüpfung und räumlichen Anordnung digitaler Baukomponenten; technische Grundlage des Entwurfswerkzeugs.}
\GlossaryDefine{Human-Interface-Design}{Gestaltung der Bedienoberflächen und Interaktionsabläufe zwischen Nutzenden und der digitalen Arbeitsumgebung.}
\GlossaryDefine{User-Experience}{Wahrnehmung und Einbindung der digitalen Arbeitsumgebung in den tatsächlichen Arbeitsablauf der Nutzenden.}
\GlossaryDefine{Entwurfsgrammatik}{Regelwerk des Projekts dafür, wie Entwurfsvarianten gebildet, verändert und auf Zulässigkeit geprüft werden.}
\GlossaryDefine{Test-Case}{Realer oder realitätsnaher Anwendungsfall, an dem die digitale Werkzeugkette schrittweise entwickelt und geprüft wird.}

\AbbreviationDefine{AP}{Arbeitspaket}{Klar abgegrenzter Teilbereich des Projekts mit eigenen Zielen und Aufgaben.}
\AbbreviationDefine{API}{Application Programming Interface}{Festgelegte Programmierschnittstelle, über die zwei Computerprogramme Daten austauschen.}
\AbbreviationDefine{KI}{Künstliche Intelligenz}{Oberbegriff für Computerprogramme, die Aufgaben lösen, für die normalerweise menschliches Denken nötig ist.}
\AbbreviationDefine{LLM}{Large-Language-Model}{Mit sehr großen Textmengen trainiertes Computerprogramm, das Texte verstehen, beantworten und formulieren kann.}
\AbbreviationDefine{HID}{Human Interface Design}{Gestaltung, wie Menschen ein Computerprogramm bedienen und mit ihm interagieren.}
\AbbreviationDefine{REST}{Representational State Transfer}{Verbreitetes technisches Verfahren, mit dem Computerprogramme über das Internet standardisiert Daten austauschen.}
\AbbreviationDefine{NGS}{Nachhaltige Gebäudesysteme}{Am Projekt beteiligte Abteilung der Leibniz Universität Hannover.}
\AbbreviationDefine{KET}{Konstruktives Entwerfen und Tragwerksplanung}{Am Projekt beteiligte Abteilung der Universität der Künste Berlin.}
\AbbreviationDefine{BBSR}{Bundesinstitut für Bau-, Stadt- und Raumforschung}{Bundesbehörde, die dieses Forschungsvorhaben fördert.}

% Reset bold to normal (scoped to this document): the table header cells get a
% synthetic FakeBold via \SemioSansBold (defined in semio-window.sty). Redefine it
% here — after the class/window.sty have loaded — to the plain Regular face so
% header rows (Plattform, Land, …) render at normal weight without touching the
% shared style files or any other report.
\makeatletter
\renewfontfamily\SemioSansBold[Path=\semio@font@root anta/]{Anta-Regular.ttf}
\makeatother

% Zeichenvokabel für Systemkarte und Prüftabellen der Anlagen. tikz wird
% bereits von semio-window.sty geladen; hier nur zusätzliche Bibliotheken und
% die Stile. Identisch zum Zwischenbericht, damit die Anlagen unverändert
% übernommen werden können.
\usetikzlibrary{matrix,positioning,fit,calc}
\makeatletter
\tikzset{
  semio/.style={every node/.append style={outer sep=0pt}},
  semio cell/.style={
    draw=semio-chrome-border-normal, line width=\semio@stroke@hairline,
    fill=semio-chrome-canvas, align=left, anchor=center,
    inner xsep=0.55em, inner ysep=0.34em,
    font=\SemioSans\fontsize{7.2pt}{9.4pt}\selectfont,
    text=semio-chrome-text-normal},
  semio module/.style={semio cell, fill=semio-chrome-window, text=semio-chrome-foreground},
  semio band/.style={draw=semio-chrome-border-emphasized,
    line width=\semio@stroke@default, fill=none},
  semio chip/.style={draw=semio-chrome-border-normal,
    line width=\semio@stroke@hairline, fill=semio-chrome-panel,
    text=semio-chrome-foreground,
    font=\SemioMono\fontsize{7.2pt}{8.4pt}\selectfont,
    inner xsep=0.4em, inner ysep=0.2em},
  semio edge/.style={draw=semio-chrome-border-emphasized,
    line width=\semio@stroke@hairline, -latex},
  semio edge synth/.style={semio edge, dash pattern=on 2.25pt off 1.5pt},
  semio bar/.style={draw=semio-chrome-border-normal,
    line width=\semio@stroke@hairline, fill=semio-chrome-panel},
}
\makeatother
% Hürden-Baum (Anlage HW): eigener Mechanismus semio-tree.sty (SemioTree /
% SemioTreeGroup / SemioTreeLeaf), von semio.cls geladen wie semio-graph.sty --
% keine Diagramm-Vokabel lebt im Dokument-Preamble.

\begin{document}

\makecoverpages

\maketableofcontents

% Anlage 8 verlangt den Förderhinweis im vorgeschriebenen Wortlaut, prominent
% platziert, deutsch und englisch.
\makefundingacknowledgement

\makemainmatter

% Gliederung wörtlich nach Anlage 8 „Berichtswesen", Version 2.3 vom 12.02.2025,
% Abschnitt „Forschungsbericht". Jedes Kapitel trägt vorerst den BBSR-Auftrag
% als Platzhalter; der Fließtext wird schrittweise befüllt.

\chapter*{Kurzbeschreibung und Abstract}

Entwerfen mit wiederverwendeten Baukomponenten ist ästhetisch, historisch und ökologisch reizvoll, aber entwurfstechnisch anspruchsvoll und komplex, da funktionale, konstruktive und qualitative Abhängigkeiten direkt beim Entwerfen berücksichtigt werden müssen.
Die im Projekt entwickelte End-To-End-Plattform, welche gleichzeitig Katalog-, Entwurfs- und Leistungsbewertungs-Tool beinhaltet, ermöglicht es, in dieser Situation effektiv Lösungen zu entwickeln.
Das Performance-Assessment hinsichtlich Energie und Tragwerk erlaubt frühe multidisziplinäre Optimierung.
Mithilfe einer KI-Assistenz können auch große Bestandspools, in welchen es explosionsartig viele Kombinationsmöglichkeiten gibt, im Entwurfsprozess berücksichtigt werden.
Die Assistenz schlägt für die jeweilige Entwurfssituation passende Baukomponenten und Verknüpfungen vor.
Basierend auf einer semantischen Repräsentation der Baukomponenten und des komponentenbasierten Entwurfes wird die generative Stärke der vortrainierten Large-Language-Models direkt für den Entwurf genutzt.
Über eine Framework-unabhängige Schnittstellendefinition wird es bestehenden Bauteilbörsen ermöglicht, ihre Baukomponenten nahtlos maschinenlesbar zu veröffentlichen, sodass die Integrierbarkeit in die Plattform gewährleistet ist.
Das Schema dieser Schnittstellendefinition beinhaltet neben Fotos und den üblichen Metadaten wie Ort und Preis auch 3D-Modelle und Performance-Kennwerte wie das Treibhauspotential, welche direkt von den Tools verwendet werden.
Ein Modellierungstool, welches moderne Bild-zu-3D-KI-Modelle benutzt, reduziert den Aufbereitungsaufwand der Bauteilbörsen.
Durch die open-source Veröffentlichung der gesamten Plattform, deren Dokumentation und den Aufbau einer Community werden die Langlebigkeit und das Vertrauen gestärkt.
Das Self-Hosting ermöglicht einerseits eine Breitenwirkung, wodurch der zukünftige Unterhalt und neue Featureentwicklung effizient geteilt wird.
Andererseits bleibt die Datensouveränität gewährleistet.

Englischer Abstract.

\chapter{Einführung}
\label{kapitel:einfuehrung}
Das Themenfeld und den Untersuchungsgegenstand darstellen: Was wurde beforscht?

\chapter{Problemstellung}
\label{kapitel:problemstellung}
Stand der Forschung und der Baupraxis, darauf aufbauend die Forschungslücke und den Entwicklungsbedarf aufzeigen: Warum wurde geforscht?

\chapter{Zielstellung}
\label{kapitel:zielstellung}
Die konkreten Projektziele sowie die übergeordneten Ziele und den Beitrag des Projekts dazu benennen: Wofür wurde geforscht?

\chapter{Forschungsdesign}
\label{kapitel:forschungsdesign}
Arbeitshypothesen, methodischen Ansatz, Projektteam, Projektorganisation und Kooperationspartner sowie Arbeitsschritte und Meilensteine beschreiben: Wie wurde geforscht?

\makeworkpackages{%
  \SemioTableTwo{%
    \SemioTableHeaderRow{Arbeitspaket & Federführung}
    \SemioTableRow{AP-Erfahrung & UdK Berlin \SemioSlash{} KET}
    \SemioTableRow{AP-Plattform & Leibniz Universität Hannover \SemioSlash{} NGS}
    \SemioTableRow{AP-Tool & Leibniz Universität Hannover \SemioSlash{} NGS}
    \SemioTableRow{AP-Validierung & UdK Berlin \SemioSlash{} KET}
  }%
}

\chapter{Projektverlauf}
\label{kapitel:projektverlauf}
Hauptteil des Berichts als Beweisführung: durchgeführte Arbeiten und Erkenntnisse einschließlich Zwischenergebnissen und Meilensteinen darstellen, Änderungen gegenüber dem ursprünglichen Antrag beschreiben und begründen.

\chapter{Ergebnisse}
\label{kapitel:ergebnisse}
Zwischenergebnisse zum Endergebnis zusammenführen und durch Bewertung, Diskussion und Zielerreichung in den wissenschaftlichen Kontext einordnen; bekannt gewordene Ergebnisse Dritter einbeziehen.

\chapter{Verwertung der erzielten Ergebnisse}
\label{kapitel:verwertung}
Überblick über die Transferaktivitäten geben und den konkreten baupraktischen Anschluss des Outputs aufzeigen.

\chapter{Schlussworte}
\label{kapitel:schlussworte}
Fazit und Ausblick.

\makeappendices

% \SemioCellList — shared bullet-list cell formatter for the appendix tables.
% Duplicated from zwischenbericht.tex (not shared via a macro package): the
% Zwischenbericht still needs its own copy for Bauteilbörsen and Typologie,
% which stay there. \SemioCellListInline is not duplicated — zero uses
% repo-wide.
\ExplSyntaxOn
\ProvideDocumentCommand \SemioCellList { O{;} m }
  {
    \seq_set_split:Nnn \l_tmpa_seq {#1} {#2}
    \int_zero:N \l_tmpa_int
    \seq_map_inline:Nn \l_tmpa_seq
      {
        \int_incr:N \l_tmpa_int
        \int_compare:nNnT \l_tmpa_int > 1 { \par \noindent }
        \hangindent=0.8em \hangafter=1
        \textcolor{semio-chrome-border-normal}{\textendash}~%
        \tl_trim_spaces:n {##1}%
      }
  }
\ExplSyntaxOff

\appendixinput[P]{Projekte}{anhang/projekte}
\SemioNavShort{Eingangsmodell}% LR-navshort
\appendixinput[EM]{Eingangsmodell für Vorbemessung und Ökobilanz}{anhang/eingangsmodell}
\appendixinput[SK]{Systemkarte Skalierung der Wiederverwendung}{anhang/skalierung}
\appendixinput[NM]{Nachweismatrix dokumentierter Plattformfunktionen}{anhang/nachweismatrix}
\appendixinput[IR]{Integrationsreife der Wiederverwendungsplattformen}{anhang/integrationsreife}
\appendixinput[AN]{Akteursnetz der Bauteil-Wiederverwendung in Europa}{anhang/akteursnetz}

\appendixchapter[V]{Verzeichnisse}
\listofappendices
\listofglossaries
\listofabbreviations
\SemioReferences

\end{document}
